\documentclass[10pt,a4paper]{amsart}
\usepackage[margin=19mm]{geometry}
\usepackage{amsmath,amssymb,mathtools}
\usepackage[scr=rsfs,cal=euler]{mathalfa}
\usepackage{xfrac}
\usepackage[unicode,hidelinks,bookmarks=false]{hyperref}
\newcommand{\ie}{i.e.\ }
\newcommand{\cf}{cf.\ }
\newcommand{\resp}{respectively\ }
\newcommand{\Subsection}{\subsection}
\providecommand{\nobreakdash}{\nobreak\hbox{-}}
\setlength{\emergencystretch}{2em}
\multlinegap0pt
\numberwithin{equation}{section} 
\numberwithin{figure}{section} 
\usepackage{stmaryrd}

\newcommand{\C}{\mathbb{C}}
\newcommand{\F}{\mathbb{F}}
\newcommand{\bbH}{\mathbb{H}}
\newcommand{\M}{\mathbb{M}}
\newcommand{\N}{\mathbb{N}}
\newcommand{\PP}{\mathbb{P}}
\newcommand{\Q}{\mathbb{Q}}
\newcommand{\R}{\mathbb{R}}
\newcommand{\RR}{\mathbb{R}}
\newcommand{\Z}{\mathbb{Z}}
\newcommand{\ZZ}{\mathbb{Z}}
\newcommand{\cF}{\mathcal{F}}
\newcommand{\cK}{\mathcal{K}}
\newcommand{\cL}{\mathcal{L}}
\newcommand{\cW}{\mathcal{W}}

\DeclareMathOperator{\Des}{Des}
\DeclareMathOperator{\des}{des}
\DeclareMathOperator{\desc}{desc}
\DeclareMathOperator{\Span}{Span}
\DeclareMathOperator{\OI}{OI}
\DeclareMathOperator{\ann}{ann}
\DeclareMathOperator{\Supp}{Supp}
\DeclareMathOperator{\PG}{PG}
\DeclareMathOperator{\mult}{mult}
\DeclareMathOperator{\wt}{wt}
\DeclareMathOperator{\topp}{top}
\DeclareMathOperator{\rk}{rk}

\newcommand*\ps[1]{{\llbracket#1\rrbracket}}

\newtheorem{theorem}{Theorem}[section]
\newtheorem{lemma}[theorem]{Lemma}
\newtheorem{prop}[theorem]{Proposition}
\newtheorem{coro}[theorem]{Corollary}
\theoremstyle{definition}
\newtheorem{defi}[theorem]{Definition}
\theoremstyle{remark}
\newtheorem{rema}[theorem]{Remark}
\newtheorem{exam}[theorem]{Example}

\title[Postnikov tree expansion]{Matroidal mixed Eulerian numbers: original Theorem 5.2}
\author{Eric Katz and Max Kutler}
\date{Source: Algebraic Combinatorics 7 (2024), publisher version; DOI 10.5802/alco.382}
\begin{document}\maketitle
\noindent\textit{Editorial scope.} The counted unit is the original published Theorem 5.2, the flat-filled increasing binary-tree expansion of an arbitrary product of hypersimplex classes in $A^*(M)$, together with its complete author proof, both reproduced exactly as printed. Retained as uncounted necessary core: the section 2 matroid Chow-ring notation (closure, $A^*(M)$, the generators $x_F$, the degree map, the $\theta_i$ pullback), the hypersimplex classes $\gamma_k$, the over-intersection $\OI$ and multiplicity $\mult$ definitions, Definition 2.1, Lemma 2.2 with its complete proof, the $\psi_{\cF}$ flag decomposition, and Lemma 2.6 with its complete proof, which is the source-local flat-intersection multiplication rule invoked in the counted proof; from section 5 the binary search order, increasing labelings, flat-fillings, $v$-compatibility, one-vertex extensions and both weight functions. Sections 3, 4, 6 and 7, Theorem 2.5, Corollary 2.7, the illustrative Examples 5.1 and 5.4, Corollary 5.3 and the two purely illustrative TikZ figures are not selected and are not used to inflate the difficulty. Original printed slips are preserved literally without mathematical repair: the doubled word in the proof of Lemma 2.2, and the two occurrences of $F(\ell(a))$ inside the definition of the weights where the surrounding product is indexed by $b$ and the text refers to $|F(r(b))|$.
\setcounter{section}{1}
\section{Matroid Chow rings and hypersimplex classes} \label{s:matroids}

\subsection{Matroids and matroid Chow rings}

We begin by reviewing some notions of matroids following \cite{AHK}.
Let $E$ denote the set $\{0,1,\dots,n\}$. Let $S_{E}$ denote the group of bijections from $E$ to itself.
Usually, we will take $M$ to be a rank $r+1$ matroid on $E$ given by a rank function $\rk\colon 2^{E}\to\Z_{\geq 0}$. 
We will denote the  {\em Boolean matroid} on $E$ by $U_{n+1,n+1}$; it is characterized by $\rk(I)=|I|$, that is, every subset of $E$ is a flat.
For $I\subseteq E$, let $\overline{I}$ denote the closure of $I$ with respect to $M$.
All matroids will be loopless unless otherwise noted. For a flat $F$, let $M_F$ denote the contraction of $M$ at $F$, i.e. the matroid whose underlying lattice of flats is the interval $[F,\hat{1}]$. Let $M^F$ denote the restriction to $F$, which has lattice $[\hat{0},F]$

In $\R^{E}$, let $e_0,\dots,e_n$ denote the standard unit basis vectors. Write $\mathbf{1}=e_0+\dots+e_n$.
Set $N_{\R}=\R^{E}/\R\mathbf{1}$ where we will conflate the $e_i$'s with their images in $N_{\R}$. For a subset $S\subset E$, write
\[e_S=\sum_{i\in S} e_i \in N_{\R}.\]
For a chain of subsets 
\[\mathcal{S}=\{\varnothing\subsetneq S_1\subsetneq \dots\subsetneq S_k\subsetneq E\},\]
write $\sigma_{\mathcal{S}}\subset N_{\R}$ for the cone
\[\sigma_{\mathcal{S}}=\Span_{\geq 0}(e_{S_1},\dots,e_{S_k}).\]
The permutohedral fan $\Delta_{E}$ is the fan in $N_{\R}$ whose cones are $\sigma_{\mathcal{S}}$ as $\mathcal{S}$ ranges over all chains of subsets.
Attached to $\Delta_{E}$ is the permutohedral toric variety $X(\Delta_{E})$.
For $w\in S_{E}$, let $\sigma_w$ be the cone in $\Delta_{E}$ given by
\[\Span_{\geq 0}(e_{S_1},\dots,e_{S_n}),\]
where $S_1\subsetneq \dots\subsetneq S_n$ is the chain of subsets
\[\{w(0)\}\subsetneq \{w(0),w(1)\}\subsetneq\dots\subsetneq \{w(0),w(1),\dots,w(n)\}.\]

The matroid Chow ring \cite{FY:chow}, $A^*(M)=\Z[x_F]/(I+J)$, is generated by $x_F$ for proper non-empty flats $F$ with relations
\begin{align*}
I&=\left\langle x_{F_1}x_{F_2}\mid F_1,F_2 \text{ are not comparable}\right\rangle\\
J&=\left\langle \sum_{F\ni i} x_F- \sum_{F\ni j} x_F\mid i,j\in E\right\rangle.
\end{align*}
There is a natural homomorphism 
\[A^*(U_{n+1,n+1})\to A^*(M)\]
given by
\[x_S\mapsto 
\begin{cases}
x_S &\text{if $S$ is a flat of $M$}\\
0 &\text{else}.
\end{cases}\]
We will sometimes write $x_S$ for the image of $x_S$ under this homomorphism even when $S$ is not a flat of $M$. For a chain of subsets 
\[\mathcal{S}=\{\varnothing\subsetneq S_1\subsetneq\dots\subsetneq S_k\subsetneq E\},\]
we will write $x_{\mathcal{S}}\coloneqq x_{S_1}\dots x_{S_k}$.
By \cite{AHK}, $A^*(M)$ is a Poincar\'{e} duality ring of dimension $r$ equipped with a degree map
\[\deg_M\colon A^r(M)\to \Z,\]
characterized by the property that for any full flag of flats
\[\{\varnothing \subsetneq F_1\subsetneq F_2\subsetneq \dots\subsetneq F_r\subsetneq E\},\]
we have $\deg(x_{F_1}x_{F_2}\dots x_{F_r})=1.$

Write $\theta_i\colon A^*(M\setminus i)\to A^*(M)$ for the pullback defined in \cite[Section~3]{BHMPW:semismall} as
\[\theta_i(x_F)=x_F+x_{F\cup i}.\]
If $i$ is not a coloop of $M$, then $\theta_i$ commutes with the degree map in the sense
\[\deg_M\circ \theta_i=\deg_{M\setminus i}.\]
If $i$ is a coloop, for dimensional reasons, $\deg_M\circ\theta_i=0$. Henceforth, we will write $A^*(M)$ for $A^*(M)\otimes\R$.

\subsection{Hypersimplex classes}

Let $S_{E}$ act on the standard unit basis vectors in $\RR^{E}$. The standard $k$-hypersimplex, $\Delta(n+1,k)$, is the convex hull of all vectors in the orbit of $e_{\{0,\dots,k-1\}}$ under this action. It lies in the hyperplane $x_0+x_1+\dots+x_n=k$. 
Attached to it is a class 
\[\gamma_{n+1-k}\in A^1(X(\Delta_{E}))\]
in the Chow cohomology ring of the permutohedral toric variety. It arises as the non-equivariant restriction of the support function to any translate $\Delta(n+1,k)-w$ for a vector $w$ with $w_0+\dots+w_n=k$. By \cite[Section~6.1]{CoxLittleSchenck}, this corresponds to the class in the equivariant Chow ring $A_T^1(\Delta_{E})$ given by the support function on $N_{\R}$,
\[\varphi(u)=\min_{v\in \Delta(n,k)-w}(u\cdot v)\\
=\min_{|T|=k}(u\cdot e_T)-u\cdot w.
\]
The corresponding non-equivariant class is obtained as
\[-\sum_S \varphi(e_S)x_S.\]
We use the convention that $\gamma_k=0$ for $k\leq 0$ or $k\geq n+1$.

\begin{defi}
For a vector $(c_1,\dots,c_n)$ of nonnegative integers with $c_1+\dots+c_n=r$, we define {\em the matroidal mixed Eulerian number}
\[A_{c_1,\dots,c_n}(M)=\deg_M(\gamma_1^{c_1}\dots\gamma_n^{c_n}).\]
\end{defi}

Because the $\gamma_k$'s are given by convex functions and hence are nef, $A_{c_1,\dots,c_n}(M)$ is always nonnegative. For $M=U_{n+1,n+1}$, these specialize to Postnikov's mixed Eulerian numbers. Indeed, those numbers are described as mixed volumes of hypersimplexes. The matroidal mixed Eulerian numbers specialize to intersection numbers in $A^*(U_{n+1,n+1})=A^*(X(\Delta_{E}))$. The connection then follows from the relationship between toric intersection theory and the polytope algebra \cite{FultonSturmfels}.

For a finite set $U$ and
$S,T\subseteq U$, the {\em over-intersection} of $S$ and $T$ in $U$ is
\[\OI_{U}(S,T)=|S\cap T|-\max(0,|S|+|T|-|U|),\]
i.e. the quantity by which the size of $S\cap T$ exceeds that which is expected for a ``generic'' choice of $S$ and $T$.
Let $\mult_{U}(|S|,k)=\min(|S|,k)-\frac{k}{|U|}|S|$.

\begin{lemma} \label{l:weights} We have the following identities for $\gamma_k$ (as elements of $A^1(M))$:
\begin{enumerate}
    \item $\gamma_k=\sum_{S\subset E} \OI_{E}(S,T)x_S$ for any set $T$ with $|T|=n+1-k$, and
    \item $\gamma_k=\sum_{S\subset E} \mult_{E}(|S|,k) x_S$.
\end{enumerate}
\end{lemma}

\begin{proof}
For the first identity, take $w=e_T$. Observe that that $e_S\cdot w=|S\cap T|$ and 
\[\min_{|U|=n-k+1}(e_S\cdot e_U)=\max(0,|S|+(n-k+1)-n).\]
Consequently, 
\[-\sum_S \varphi(e_S)x_S=\OI_{E}(S,T)x_S.\]

For the second identity, take $w=\frac{n+1-k}{n+1}e_{E}$, and note that $e_S\cdot w=\frac{n+1-k}{n+1}|S|$. 
\end{proof}

For a flag of flats 
\[\varnothing\subsetneq F_1\subsetneq F_2\subsetneq\dots\subsetneq F_c\subsetneq E,\]
by inducting on \cite[Prop~2.25]{BHMPW:semismall}, there is an isomorphism $\psi_{\cF}$,
\[
A^*(M)/\ann(x_{F_1}x_{F_2}\dots x_{F_c})\cong
A^*(M^{F_1})\otimes A^*(M_{F_1}^{F_2})\otimes\dots \otimes A^*(M_{F_{c-a}}^{F_c})\otimes A^*(M_{F_c})
\]
induced by multiplication by $x_{F_1}x_{F_2}\dots x_{F_c}$.
Here, $M_{F_j}^{F_{j+1}}$ is a matroid on $F_{j+1}\setminus F_j$ with lattice of flats $[F_j,F_{j+1}]$.
Moreover, if one equips $A^*(M)/\ann(x_{F_1}x_{F_2}\dots x_{F_c})$ with the degree map
\[\deg\colon A^{r-c}(M)/\ann(x_{F_1}x_{F_2}\dots x_{F_c})\to\Z,\ y\mapsto \deg_M(yx_{F_1}x_{F_2}\dots x_{F_c}),\]
and equips $A^*(M^{F_1})\otimes A^*(M_{F_1}^{F_2})\otimes\dots \otimes A^*(M_{F_{c-a}}^{F_c})\otimes A^*(M_{F_c})$
with 
\[\deg_{M^{F_1}}\otimes\deg_{M_{F_1}^{F_2}}\otimes \dots \otimes \deg_{M_{F_c}},\]
then $\psi_{\cF}$ commutes with degree maps.

We now compute the images $\psi_{\cF}(\gamma_k)$ following \cite[Lemma 6.5]{BST}.

\setcounter{theorem}{5}
\begin{lemma} \label{l:productwithflag}
The image of $\gamma_k$ under $\psi_{\cF}$ is 
\[\begin{cases}
0 &\text{if $k=|F_j|$ for some $j$}\\
1^{\otimes j}\otimes \gamma_{k-|F_j|}\otimes 1^{\otimes(c-j)}&\text{if $|F_j|<k<|F_{j+1}|$ for some $j$}.
\end{cases}
\]
\end{lemma}

\begin{proof}
There is a unique $j$ such that $|F_j|\leq k<|F_{j+1}|$. Pick $T'\subset F_{j+1}\setminus F_j$ of size $|F_{j+1}|- k$, and set $T=(E\setminus F_{j+1})\cup T'$. Then,
  \begin{align*}
      x_{F_1}x_{F_2}\dots x_{F_c}\gamma_k
      &=  x_{F_1}x_{F_2}\dots x_{F_c}\left(\sum_F \OI_{E}(F,T)x_F\right)\\
      &=  x_{F_1}x_{F_2}\dots x_{F_c}\left(\sum_{F\colon F_j\subseteq F\subset F_{j+1}} \OI_{E}(F,T)x_F\right)\\
      &=  x_{F_1}x_{F_2}\dots x_{F_c}\left(\sum_{F\colon F_j\subseteq F\subset F_{j+1}} \OI_{F_{j+1}\setminus F_j}(F\setminus F_j,T')x_F\right)
  \end{align*}
  where the second equality is a consequence of our choice of $T$.
  The conclusion follows from noting that the sum is a description of $\gamma_{k-|F_j|}$ in $A^*(M_{F_j}^{F_{j+1}})$.
\end{proof}

\setcounter{section}{4}
\section{Postnikov trees} \label{s:trees}

In \cite[Section 17]{Postnikov:permutohedra}, Postnikov gave a combinatorial interpretation of mixed Eulerian numbers as weighted counts of certain binary trees. We extend this construction to express an arbitrary product of the $\gamma_i$'s in $A^*(M)$ as a weighted sum of monomials $x_{\cF}$, for $\cF$ a flag of flats of $M$.

Let $T$ be a finite binary tree. 
The {\em binary search order} on the vertices of $T$ is the transitive closure of the relations
\begin{enumerate}
    \item $b\in L_a$ implies $b<a$ and
    \item $b\in R_a$ implies $a<b$,
\end{enumerate}
where $L_a$ and $R_a$ denote the left and right branches, respectively, under $a$. 
Let $\desc(a,T)\coloneqq L_a\cup\{a\}\cup R_a$ be the set of all descendants of $a$.

An \emph{increasing labeling} of the vertex set of $T$ is a bijection
\[\sigma\colon V(T) \to \{1,\dots,k\}\]
such that $\sigma(b)\geq \sigma(a)$ whenever $b\in\desc(a,T)$.
An \emph{increasing binary tree} is a pair $(T,\sigma)$ where $T$ is a binary tree with increasing labeling of $\sigma$. It is well-known that increasing binary trees on the vertex set $\{1, \ldots, k\}$ are in bijection with permutations of $\{1, \ldots, k\}$ \cite[Section 1.5]{Stanley}.

Let $\cL(M)$ denote the lattice of flats of $M$. A {\em flat-filling} of $T$ is a function
\[F\colon V(T)\to \cL(M)\setminus\{\varnothing, E\}\]
such that $a<b$ implies $F(a)\subsetneq F(b)$. Consequently, the image of a flat-filling must be a $k$-step flag of flats $\cF(T,F)$.

A {\em flat-filled increasing binary tree} is a triple $(T,\sigma,F)$, where $(T,\sigma)$ is an increasing binary tree and $F\colon V(T)\to \cL(M)$ 
 is a flat-filling. We will give a necessary {\em compatibility condition} for $(T,\sigma,F)$ to contribute a multiple of $x_{\cF(T,F)}$ in a particular monomial expansion of $\gamma_{v_1}\dots\gamma_{v_k}$ where $k=|V(T)|$. This will allow us to construct a flat-filled increasing binary tree vertex-by-vertex in an order determined by $\sigma$. 
For $i \in \{1,\ldots,k\}$, write $T_{\leq i}$ for the subgraph of $T$ induced by the vertex set $\sigma^{-1}(\{1,\ldots,i\})$. We will also write $F$ for the flat-filling on $T_{\leq i}$ given by restricting $F$ from $T$.
Observe that $T_{\leq i}$ is also a binary tree, and the binary search order on $T_{\leq i}$ is the restriction of the binary search order on $T$. For a vertex $b$, let $\ell(b)$ and $r(b)$ denote $b$'s immediate predecessor and successor, respectively, in the binary search order on $T_{\leq \sigma(b)}$.

For a vector $v = (v_1, \ldots, v_k)$, we define a flat-filled increasing binary tree $(T,\sigma,\cF)$ to be {\em $v$-compatible} if for all $b\in V(T)$
\[ |F(\ell(b))| < v_{\sigma(b)} < |F(r(b))|\]
where if $b$ is the minimal (resp. maximal) element in the binary search order on $T$, we set $F(\ell(b))=\varnothing$ (resp. $F(r(b))=E$).
In light of Lemma~\ref{l:productwithflag} and the discussion above, this will translate to the condition that $\gamma_{v_{\sigma(b)}}$ could have added the flat $F(b)$ to $\cF(T_{\leq \sigma(b)-1},F)$ to create $\cF(T_{\leq \sigma(b)},F)$.

We define a {\em one-vertex extension} of a flat-filled increasing tree $(T,\sigma,F)$ on $k$ vertices to be a flat-filled increasing tree $(T',\sigma',F')$ on $k+1$ vertices such that $T'_{\leq k} = T$, $\sigma'|_T = \sigma$, and $F'|_{V(T)} = F$.
If the new vertex is called $b$, then $F'(\ell(b))\subsetneq F'(b)\subsetneq F'(r(b))$.


We can choose between two possible natural weights for a $v$-compatible flat-filled increasing binary tree $(T,\sigma,F)$:
\[\wt^{\OI}_v(T,\sigma,F) \coloneqq \prod_{b\in V(T)} \OI_{F(r(b))\setminus F(\ell(b))} (F(b) \setminus F(\ell(b)), U_b)\]
where $U_b$ is the set of the largest $|F(r(b))| - v_{\sigma(b)}$ elements of $F(r(b)) \setminus F(\ell(a))$, or
\[\wt^{\mult}_v(T,\sigma,F) \coloneqq \prod_{b\in V(T)} \mult_{F(r(b))\setminus F(\ell(b))} (F(b) \setminus F(\ell(b)), v_{\sigma(b)}-|F(\ell(a))|).\]
We set the weight of the empty binary tree to be $1$.

\setcounter{theorem}{1}
\begin{theorem} \label{t:postnikov}
    Let $v = (v_1, \ldots, v_k) \in \N^k$. Then we have the following equality in $A^*(M)$:
    \[ \gamma_{v_1} \cdots \gamma_{v_k} = \sum_{(T,\sigma,F)} \wt_v(T,\sigma,F) \, x_{\cF(T,F)}, \]
    (for either set of weights) where the sum is over all $v$-compatible flat-filled increasing binary trees on $k$ vertices.
\end{theorem}

\begin{proof}
    We will state the proof for the weight function $\wt^{\OI}_v$. The proof using $\wt^{\mult}_v$ is identical.
    The proof is by induction on $k$. For $k=0$, it is trivially true.

    We first claim that for any flat-filled increasing binary tree $(T,\sigma,F)$ with $k$ vertices,
    \[\gamma_{v_{k+1}}x_{\cF(T,F)}=\sum_{(T',\sigma',F')} \OI_{F'(r(b))\setminus F'(\ell(b))} (F'(b) \setminus F'(\ell(b)), U_b)x_{\cF(T',F')}
    \]
    where the sum is over one-vertex extensions of $(T,\sigma,F)$ by a vertex $b$ for which $|F'(\ell(b))|<v_{k+1}<|F'(r(b))|$. 
    Write 
    \[\cF(T,F)=\{\varnothing=F_0\subsetneq F_1\subsetneq F_2\subsetneq \dots\subsetneq F_k\subsetneq E\}.\]
    If $v_{k+1}=|F_j|$ for any $j$, then $\gamma_{v_{k+1}}x_{\cF(T',F')}=0$ and there are no one-vertex extensions $(T',\sigma',F')$ with $|F'(\ell(b))|<v_{k+1}<|F'(r(b))|$.
    Otherwise,  if $|F_j|<v_{k+1}<|F_{j+1}|$, by Lemma~\ref{l:productwithflag}, 
    \[\gamma_vx_{\cF(T,F)}=\sum_{G \colon F_j \subsetneq G \subsetneq F_{j+1}}\OI_{F_{j+1}\setminus F_j} (G \setminus F_j, U)x_Gx_{\cF(T',F')}
    \]
    where $U$ denotes the set of the largest $|F_{j+1}|-v_{k+1}$ elements of $F_{j+1}\setminus F_j$.
    To each $G$ occurring in the above sum, we produce a one-vertex extension of $(T,\sigma,F)$ by adjoining a vertex $b$ such that $F^{-1}(F_j)<b<F^{-1}(F_{j+1})$ in the binary search order and setting $\sigma(b)=k+1$, $F(b)=G$. All one-vertex extensions by $b$ for which $|F'(\ell(b))|<v_{k+1}<|F'(r(b))|$ arise in this fashion.

    Now, we give the inductive step. For $v'=(v_1,\dots,v_{k+1})$, set $v=(v_1,\dots,v_k)$. Write
    \[\gamma_{v_1}\dots\gamma_{v_k}=\sum_{(T,\sigma,F)} \wt^{\OI}_v(T,\sigma,F) \, x_{\cF(T,F)}.
    \]
    Multiply both sides by $\gamma_{v_{k+1}}$. Each choice of $(T,\sigma,F)$ contributes a sum over one-vertex extensions $(T',\sigma',F')$ with $|F'(\ell(b))|<v_{k+1}<|F'(r(b))|$ weighted by an over-intersection term. The condition on the size of flats adjacent to $b$ is exactly $v$-compatibility. Because the product of $\wt^{\OI}_v(T,\sigma,F)$ with the over-intersection term is exactly $\wt^{\OI}_{v'}(T',\sigma',F')$, the conclusion follows.
 
\end{proof}

\begin{thebibliography}{99}
\bibitem[1]{AHK} Karim Adiprasito, June Huh, and Eric Katz, \emph{Hodge theory for combinatorial geometries}, Ann. of Math. (2) \textbf{188} (2018), no. 2, 381--452.
\bibitem[3]{BST} Andrew Berget, Hunter Spink, and Dennis Tseng, \emph{Log-concavity of matroid $h$-vectors and mixed Eulerian numbers}, Duke Math. J. \textbf{172} (2023), no. 18, 3475--3520.
\bibitem[5]{BHMPW:semismall} Tom Braden, June Huh, Jacob P. Matherne, Nicholas Proudfoot, and Botong Wang, \emph{A semismall decomposition of the Chow ring of a matroid}, Adv. Math. \textbf{409} (2022), article no. 108646 (49 pages).
\bibitem[8]{CoxLittleSchenck} David A. Cox, John B. Little, and Henry K. Schenck, \emph{Toric varieties}, Graduate Studies in Mathematics, vol. 124, American Mathematical Society, Providence, RI, 2011.
\bibitem[15]{FY:chow} Eva Maria Feichtner and Sergey Yuzvinsky, \emph{Chow rings of toric varieties defined by atomic lattices}, Invent. Math. \textbf{155} (2004), no. 3, 515--536.
\bibitem[16]{FultonSturmfels} William Fulton and Bernd Sturmfels, \emph{Intersection theory on toric varieties}, Topology \textbf{36} (1997), no. 2, 335--353.
\bibitem[23]{Postnikov:permutohedra} Alexander Postnikov, \emph{Permutohedra, associahedra, and beyond}, Int. Math. Res. Not. IMRN (2009), no. 6, 1026--1106.
\bibitem[24]{Stanley} Richard P. Stanley, \emph{Enumerative combinatorics. Volume 1}, second ed., Cambridge Studies in Advanced Mathematics, vol. 49, Cambridge University Press, Cambridge, 2012.
\end{thebibliography}
\end{document}
